\section*{Author source: syntax and inference rules (30:3--4)}
Let $\mathsf{Var}$ be a fixed countable set of propositional variables. The set $\mathsf{Fm}$ of formulas is defined by the following grammar:
\[
 A,B::=p\mid1\mid0\mid A\wedge B\mid A\vee B\mid A\cdot B\mid A\to B\quad(p\in\mathsf{Var}).
\]
Here we adopt the standard notation in substructural logics. In linear logic, the connectives $1,0,\wedge,\vee,\cdot,\to$ are respectively written as $1,\bot,\&,\oplus,\otimes,\multimap$. We write $\neg A$ for $A\to0$, and $A\leftrightarrow B$ for $(A\to B)\wedge(B\to A)$.
A sequent is of the form $\Gamma\Rightarrow\Pi$, where $\Gamma$ is a multiset of formulas and $\Pi$ is either a singleton of a formula or an empty set. Inference rules of the sequent calculus for $\FLew$ are given in Figure 1.
\[
 \frac{}{A\Rightarrow A}\ (\mathrm{init})\qquad
 \frac{\Gamma\Rightarrow A\quad A,\Delta\Rightarrow\Pi}{\Gamma,\Delta\Rightarrow\Pi}\ (\mathrm{cut})
\]
\[
 \frac{\Delta\Rightarrow\Pi}{\Gamma,\Delta\Rightarrow\Pi}\ (i)\qquad
 \frac{\Gamma\Rightarrow}{\Gamma,\Delta\Rightarrow\Pi}\ (o)
\]
\[
 \frac{A_i,\Delta\Rightarrow\Pi}{A_1\wedge A_2,\Delta\Rightarrow\Pi}\ (\wedge\ell)\qquad
 \frac{\Gamma\Rightarrow A_1\quad\Gamma\Rightarrow A_2}{\Gamma\Rightarrow A_1\cdot A_2}\ (\wedge r)
\]
\[
 \frac{\Delta\Rightarrow\Pi}{1,\Delta\Rightarrow\Pi}\ (1\ell)\qquad \frac{}{\Rightarrow1}\ (1r)
\]
\[
 \frac{A_1,\Delta\Rightarrow\Pi\quad A_2,\Delta\Rightarrow\Pi}{A_1\vee A_2,\Delta\Rightarrow\Pi}\ (\vee\ell)\qquad
 \frac{\Gamma\Rightarrow A_i}{\Gamma\Rightarrow A_1\vee A_2}\ (\vee r)
\]
\[
 \frac{}{0\Rightarrow}\ (0\ell)\qquad \frac{\Gamma\Rightarrow}{\Gamma\Rightarrow0}\ (0r)
\]
\[
 \frac{A,B,\Delta\Rightarrow\Pi}{A\cdot B,\Delta\Rightarrow\Pi}\ (\cdot\ell)\qquad
 \frac{\Gamma\Rightarrow A\quad\Delta\Rightarrow B}{\Gamma,\Delta\Rightarrow A\cdot B}\ (\cdot r)
\]
\[
 \frac{\Gamma\Rightarrow A\quad B,\Delta\Rightarrow\Pi}{\Gamma,A\to B,\Delta\Rightarrow\Pi}\ (\to\ell)\qquad
 \frac{\Gamma,A\Rightarrow B}{\Gamma\Rightarrow A\to B}\ (\to r)
\]
\begin{center}Figure 1. Sequent calculus for $\FLew$.\end{center}
Given a formula $A\in\mathsf{Fm}$, we write $A^{\mathrm{seq}}$ for the sequent $\Rightarrow A$ with empty antecedent.
Given a set $\Phi\cup\{A\}$ of formulas, we write $\Phi\vdash_{\FLew} A$ if $A^{\mathrm{seq}}$ is derivable from some sequents in $\Phi^{\mathrm{seq}}:=\{B^{\mathrm{seq}}\mid B\in\Phi\}$ by rules in Figure 1.
\paragraph{Definition 1.} A substructural logic over $\FLew$ is a set $L\subseteq\mathsf{Fm}$ which is closed under deduction in $\FLew$ and substitution:
\[
 L\vdash_{\FLew}A\Longrightarrow A\in L,\qquad
 A\in L\Longrightarrow A'\in L,\quad\text{where $A'$ is a substitution instance of $A$.}
\]
We denote by $\operatorname{Ext}(\FLew)$ the set of substructural logics over $\FLew$.
\section*{Author source: algebraic semantics (30:4--5, 30:9)}
\paragraph{Definition 2.} An $\FLew$ algebra is a tuple $P=(P,\wedge,\vee,\cdot,\to,1,0)$ such that
\begin{itemize}
\item $(P,\wedge,\vee,1,0)$ is a bounded lattice with $1$ greatest and $0$ least,
\item $(P,\cdot,1)$ is a commutative monoid, and
\item $a\cdot(-):P\to P$ is a left adjoint of $a\to(-):P\to P$ for every $a\in P$:
\[
 a\cdot x\le y\Longleftrightarrow x\le a\to y,
\]
where $x\le y$ is a shorthand for $x\vee y=y$ (or equivalently $x=x\wedge y$).
\end{itemize}
A formula $A\in\mathsf{Fm}$ can be seen as a term of $\FLew$ algebras. Hence it makes sense to consider a homomorhism $v:\mathsf{Fm}\to P$ for any $\FLew$ algebra $P$. Such $v$ is called a valuation. We say that $v$ satisfies $\Phi\subseteq\mathsf{Fm}$ if $v(A)=1$ holds for every $A\in\Phi$.
Given a class $K$ of $\FLew$ algebras and a set $\Phi\cup\{A\}$ of formulas, we write $\Phi\models_K A$ if for every $P\in K$ and every valuation $v:\mathsf{Fm}\to P$ with $P\in K$, if $v$ satisfies $\Phi$, then it also satisfies $A$. When $\Phi=\varnothing$, we write $\models_K A$. When $K=\{P\}$, we write $\Phi\models_P A$.
We have strong soundness: for any set $\Phi\cup\{A\}\subseteq\mathsf{Fm}$,
\[
 \Phi\vdash_L A\Longrightarrow\Phi\models_K A.
\]
\paragraph{Figure 3 (Łukasiewicz row).} $x*y=\max(x+y-1,0)$ and $x\to y=1-x+y$ when $x>y$. We have $x\to y=1$ when $x\le y$. Its $\FLew$ algebra is $[0,1]_{\Luk}$.
\paragraph{Theorem 10 (Completeness of major fuzzy logics; selected pair).}
Finite strong completeness holds for $(\Luk,[0,1]_{\Luk})$.
Note that implication $\to$ is continuous in $[0,1]_{\Luk}$, while it is discontinuous at $(0,0)$ in $[0,1]_G$ and $[0,1]_{\Pi}$. In fact, we have $0\to0=1$, while $x\to0=0$ for any $x>0$. Since any continuous t-norm is an ordinal sum of Gödel, Łukasiewicz and product norms [26] (see also [15,8,24]), a well-known fact is obtained.
\paragraph{Theorem 12.} Among the infinitely many t-norms, Łukasiewicz norm is the only one that makes all logical connectives $\wedge,\vee,\cdot,\to$ continuous.
\section*{Author source: homomorphism calculation used again in Theorem 31 (30:10)}
Define a (partial) valuation $v:\overline\Phi\to[0,1]_{\Luk}$ by setting $v(A_i):=c_i$. It is homomorphic for each $\star\in\{\wedge,\vee,\cdot,\to\}$. In fact, if $A_i\equiv A_j\star A_k$,
\[
v(A_i)=c_i=f_i(\mathbf c)=c_j\star c_k=v(A_j)\star v(A_k).
\]
\section*{Author source: naive set theory (30:12)}
Let $\mathsf{Var}_{\mathrm{set}}:=\{x_0,x_1,x_2,\ldots\}$ be a fixed countable set of term variables. The set $\mathsf{Tm}_{\mathrm{set}}$ of terms and the set $\mathsf{Fm}_{\mathrm{set}}$ of formulas are simultaneously defined by the following grammar:
\[
 \begin{aligned}
 t,u&::=x\mid\{x\mid A\},\\
 A,B&::=t\in u\mid1\mid0\mid A\wedge B\mid A\vee B\mid A\cdot B\mid A\to B,
 \end{aligned}
\]
where variable $x$ is bound in formula $\forall x.A$ and in term $\{x\mid A\}$. Existential quantifier is defined by $\exists y.A:=\forall x.(\forall y.(A\to x^{\bullet}))\to x^{\bullet}$, where $x$ is fresh and $x^{\bullet}:=\bullet\in x$ for a fixed closed term $\bullet$. In the sequel, we will also use notation $A\otimes B$ instead of $A\cdot B$ for readability.
The sequent calculus for $\FLew_{\mathrm{set}}$ consists of the inference rules in Figure 1 together with:
\[
 \frac{A[t/x],\Delta\Rightarrow\Pi}{t\in\{x\mid A\},\Delta\Rightarrow\Pi}\ (\in\ell)\qquad
 \frac{\Gamma\Rightarrow A[t/x]}{\Gamma\Rightarrow t\in\{x\mid A\}}\ (\in r)
\]
\[
 \frac{A[t/x],\Delta\Rightarrow\Pi}{\forall x.A,\Delta\Rightarrow\Pi}\ (\forall\ell)\qquad
 \frac{\Gamma\Rightarrow A}{\Gamma\Rightarrow\forall x.A}\ (\forall r)
\]
where rule $(\forall r)$ is subject to the eigenvariable condition.
Given a set $\Phi\cup\{A\}\subseteq\mathsf{Fm}_{\mathrm{set}}$, we write $\Phi\vdash_{\FLew_{\mathrm{set}}}A$ if $A^{\mathrm{seq}}$ is derivable from $\Phi^{\mathrm{seq}}$ by the above inference rules.
Given $L\in\operatorname{Ext}(\FLew)$, let $\widetilde L\subseteq\mathsf{Fm}_{\mathrm{set}}$ be the set of all substitution instances of propositional formulas in $L$ (in the language $\mathsf{Fm}_{\mathrm{set}}$). We define $L_{\mathrm{set}}:=\{A\in\mathsf{Fm}_{\mathrm{set}}\mid\widetilde L\vdash_{\FLew_{\mathrm{set}}}A\}$, and write $\vdash_{L_{\mathrm{set}}}A$ if $A\in L_{\mathrm{set}}$.
\section*{Author source: restricted system (30:15)}
What we will now take is the second approach. To apply the BFT argument, we have to find a suitable fragment of $\Luk_{\mathrm{set}}$, for which a hypercube of finite dimension suffices. The following fragment meets this requirement.
\paragraph{Definition 30.} System $\FLew_{\mathrm{set}}^{-}$ is obtained from $\FLew_{\mathrm{set}}$ by imposing the following restriction:
\begin{quote}
(*) when applying rule $(\forall\ell)$, $t$ must be either a variable or a closed term.
\end{quote}
We define $\Luk_{\mathrm{set}}^{-}:=\{A\in\mathsf{Fm}_{\mathrm{set}}\mid\widetilde{\Luk}\vdash_{\FLew_{\mathrm{set}}^{-}}A\}$.
\section*{Author source: complete consistency construction and proof (30:15--16)}
Let us now prove the consistency of $\Luk_{\mathrm{set}}^{-}$. We begin with some preparations. Suppose that a finite set $\Phi\subseteq\mathsf{Fm}_{\mathrm{set}}\cup\mathsf{Tm}_{\mathrm{set}}$ of formulas and terms is given. We consider its closure $\overline\Phi$ under subformulas and subterms. Explicitly, $\overline\Phi$ is the least set such that (i) $\Phi\subseteq\overline\Phi$, (ii) $A\star B\in\overline\Phi$ implies $A,B\in\overline\Phi$ for $\star\in\{\wedge,\vee,\cdot,\to\}$, (iii) $\forall x.A\in\overline\Phi$ implies $A\in\overline\Phi$, (iv) $(t\in u)\in\overline\Phi$ implies $t,u\in\overline\Phi$, and (v) $\{x\mid A\}\in\overline\Phi$ implies $A\in\overline\Phi$.
Let $\Phi_{\mathrm{ctm}}$ be the set of closed terms in $\overline\Phi$, and let
\[
 \Phi_{\mathrm{cfm}}:=\{A\sigma\in\mathsf{Fm}_{\mathrm{set}}\mid A\in\overline\Phi,\ \sigma:\mathsf{Var}_{\mathrm{set}}\to\Phi_{\mathrm{ctm}}\},
\]
where $A\sigma$ is a closed formula obtained by applying substitution $\sigma$ to $A$. As is easily observed, $\forall x.B\in\Phi_{\mathrm{cfm}}$ implies $B[t/x]\in\Phi_{\mathrm{cfm}}$ for every $t\in\Phi_{\mathrm{ctm}}$. Since $\overline\Phi$ is a finite set, so are $\Phi_{\mathrm{ctm}}$ and $\Phi_{\mathrm{cfm}}$.
Suppose that $\Phi_{\mathrm{cfm}}=\{A_1,\ldots,A_n\}$. We define a function $f=(f_1,\ldots,f_n):[0,1]^n\to[0,1]^n$ as in the proof of Theorem 14. Concretely, given $\mathbf x=(x_1,\ldots,x_n)\in[0,1]^n$, let
\[
\begin{aligned}
 f_i(\mathbf x)&:=\ell&&\text{if }A_i\equiv\ell\ (\ell=0,1),\\
 &:=x_j\star x_k&&\text{if }A_i\equiv A_j\star A_k\ (\star\in\{\wedge,\vee,\cdot,\to\}),\\
 &:=\bigwedge_{t\in\Phi_{\mathrm{ctm}}}x_{j_t}&&\substack{\text{if }A_i\equiv\forall x.B\text{ and }B[t/x]\equiv A_{j_t}\\\text{for each }t\in\Phi_{\mathrm{ctm}}\ (1\le j_t\le n)},\\
 &:=x_j&&\substack{\text{if }A_i\equiv t\in\{x\mid B\}\text{ and }B[t/x]\equiv A_j\\\text{for some }1\le j\le n},\\
 &:=0&&\text{if }A_i\equiv t\in u\text{ but the previous case does not apply.}
\end{aligned}
\]
Then $f$ is continuous by Theorem 12 and finiteness of $\Phi_{\mathrm{ctm}}$. By BFT, there is $\mathbf c=(c_1,\ldots,c_n)\in[0,1]^n$ such that $f(\mathbf c)=\mathbf c$. This gives us a function $v:\Phi_{\mathrm{cfm}}\to[0,1]_{\Luk}$ defined by $v(A_i):=c_i$.
$v$ is homomorphic with respect to propositional connectives $1,0,\wedge,\vee,\cdot,\to$ as in the proof of Theorem 14. Moreover, if $A_i\equiv t\in\{x\mid B\}$ and $A_j\equiv B[t/x]$, we have
\begin{equation}
 v(t\in\{x\mid B\})=c_i=f_i(\mathbf c)=c_j=v(B[t/x]),\tag{5}
\end{equation}
so $v$ validates rules $(\in\ell)$ and $(\in r)$ applied to formulas in $\Phi_{\mathrm{cfm}}$. Finally if $A_i\equiv\forall x.B$ and $A_{j_t}\equiv B[t/x]$ for each $t\in\Phi_{\mathrm{ctm}}$, we have
\begin{equation}
 v(\forall x.B)=c_i=f_i(\mathbf c)=\bigwedge_{t\in\Phi_{\mathrm{ctm}}}c_{j_t}
 =\bigwedge_{t\in\Phi_{\mathrm{ctm}}}v(B[t/x]).\tag{6}
\end{equation}
We are now ready to give a consistency proof.
\paragraph{Theorem 31.} $\Luk_{\mathrm{set}}^{-}$ is consistent.
\paragraph{Proof.} Suppose that $\Luk_{\mathrm{set}}^{-}$ is inconsistent. This means that there are a finite set of formulas $\Psi\subseteq\widetilde{\Luk}$ and a derivation $\pi$ of $\Rightarrow0$ from $\Psi^{\mathrm{seq}}$. Let $\Phi$ be the finite set of all formulas occurring in $\pi$, and define $\Phi_{\mathrm{ctm}},\Phi_{\mathrm{cfm}}$ and valuation $v:\Phi_{\mathrm{cfm}}\to[0,1]_{\Luk}$ as above.
We prove by induction on the derivation the following:
\begin{quote}
if sequent $\Gamma\Rightarrow\Pi$ occurs in $\pi$ and $\sigma:\mathsf{Var}_{\mathrm{set}}\to\Phi_{\mathrm{ctm}}$, then $v(\Gamma\sigma)\le v(\Pi\sigma)$.
\end{quote}
Here $v(\Gamma):=v(A_1)\cdots v(A_m)$ if $\Gamma\equiv A_1,\ldots,A_m$, and $v(\Gamma):=1$ if $\Gamma$ is empty. Likewise, $v(\Pi):=v(A)$ if $\Pi\equiv\{A\}$, and $v(\Pi):=0$ otherwise.
\begin{itemize}
\item If the derivation just consists of $\Rightarrow B$ with $B\in\Psi$, then we have $1\le v(B\sigma)$ as $B\sigma$ is a substitution instance of a propositional formula in $\Luk$.
\item If the derivation ends with a rule for a propositional connective or $(\in\ell)$ or $(\in r)$, the claim easily follows from the induction hypothesis and (5).
\item Suppose that the derivation ends with rule $(\forall r)$ inferring $\Gamma\Rightarrow\forall x.B$ from $\Gamma\Rightarrow B$. Given $\sigma:\mathsf{Var}_{\mathrm{set}}\to\Phi_{\mathrm{ctm}}$, let $\sigma'$ be a substitution that agrees with $\sigma$ on $\mathsf{Var}_{\mathrm{set}}\setminus\{x\}$ but leaves $x$ unchanged. By the induction hypothesis and the eigenvariable condition, we have $v(\Gamma\sigma)=v(\Gamma\sigma')\le v(B\sigma'[t/x])$ for every $t\in\Phi_{\mathrm{ctm}}$. Hence by (6),
\[
 v(\Gamma\sigma)\le\bigwedge_{t\in\Phi_{\mathrm{ctm}}}v(B\sigma'[t/x])
 =v(\forall x.B\sigma')=v((\forall x.B)\sigma).
\]
\item Now suppose that the derivation ends with rule $(\forall\ell)$ inferring $\forall x.B,\Gamma\Rightarrow\Pi$ from $B[t/x],\Gamma\Rightarrow\Pi$. By the induction hypothesis we have $v(B[t/x]\sigma)\cdot v(\Gamma\sigma)\le v(\Pi\sigma)$.
If $t$ is closed, then $t\sigma=t\in\Phi_{\mathrm{ctm}}$. If $t$ is a variable $y$, then $t\sigma=\sigma(y)\in\Phi_{\mathrm{ctm}}$. In any case, we have $t\sigma\in\Phi_{\mathrm{ctm}}$. So
\[
 v((\forall x.B)\sigma)=v(\forall x.B\sigma')=\bigwedge_{u\in\Phi_{\mathrm{ctm}}}v(B\sigma'[u/x])
 \le v(B\sigma'[t\sigma/x])=v(B[t/x]\sigma),
\]
where $\sigma'$ is as above. Hence we conclude $v((\forall x.B)\sigma)\cdot v(\Gamma\sigma)\le v(\Pi\sigma)$.
\end{itemize}
This is a contradiction because derivation $\pi$ ends with sequent $\Rightarrow0$.\hfill$\triangleleft$
